\section{Latent identification foundations: finite proofs}
\label{sec:annex-p16-latent-identification-foundations}

The source's foundation statements and proofs are in
\cite[Definition~3 and Lemmas~4--8, pp.~1221, 1224]{ShpitserPearl2006JointIdentification}.
The following arguments use the explicitly retained latent inputs to
prove the finite identities without an unproved latent-projection step.

\begin{annexitem}{annex:p16.latent-kernel-semantics}{thm:p16-latent-kernel-semantics}{Expanded roots, normalization and restriction after intervention}
\begin{proof}
Take the product outcome space \(\mathcal U\) with weights
\(r(u)=\prod_\alpha r_\alpha(u_\alpha)\). Finite distributivity
shows that these weights sum to one. Introduce a root vertex for
each latent input, whose equation copies its coordinate disturbance.
For each observed vertex use the stated equation, with the appropriate
root variables and observed parents as inputs, and a constant unused
disturbance. The latent coordinate disturbances and these constants
are mutually independent. Putting the latent roots first, followed
by an observed topological order, gives a topological order of the
expanded graph. All functions are total. This constructs a model of
\cref{def:p16-finite-scm}; its topological evaluation and uniqueness
are supplied by \cref{thm:p16-causal-topological-order,thm:p16-causal-solution}.

Fix \(H\), an outside assignment, and a latent tuple \(u\).
After fixing the observed equations outside \(H\), a proposed
assignment on \(H\) is the solution exactly when every retained
equation there holds. Each such assertion is the indicator in the
displayed kernel formula. Their product is therefore one exactly at
that solution and zero at all other assignments. Multiplying by the
latent tuple's mass and summing over \(u\) proves the formula.
For each fixed \(u\), summation over \(S_H\) gives one, by
existence and uniqueness of that solution. Summation over the latent
law consequently proves kernel normalization. The argument uses no
observational conditioning and includes zero-probability tuples.

Only equations indexed by \(H\) occur in this expression. The
outside observed coordinates used in them are exactly the allowed
outside parent inputs. Thus the kernel does not depend on any other
outside value.

For the restriction statement, replace the outside observed parent
arguments in each retained function by their fixed values, keep
\(\operatorname{pa}(i)\cap H\) as its observed input set, and
retain every latent input with child set \(C_\alpha\cap H\).
The latent law is unchanged. An input with no remaining child can
be omitted, since summing its mass gives one. The directed graph
among remaining observed vertices is the induced directed graph.
Two remaining vertices share an input exactly when they shared one
before, giving precisely the induced bidirected edges. This uses
the declared-input convention even if a specialized function ignores
one of its allowed arguments. The restricted equations are exactly
those just evaluated for \(Q[H]\), so they have that distribution.
No arbitrary unobserved nonroot variable has been removed or replaced.
\end{proof}
\end{annexitem}

\begin{annexitem}{annex:p16.latent-district-partition}{thm:p16-latent-district-partition}{Connectivity, spanning trees and the latent partition}
\begin{proof}
Connection by bidirected paths is reflexive through the length-zero
path and symmetric by reversing a path. Concatenating two such paths
gives a walk; deleting portions between repeated vertices gives a path.
This proves transitivity. The equivalence classes are consequently a
unique partition, and each is maximal among connected vertex sets.

For a nonempty finite connected bidirected graph, start with one
vertex. If some vertex remains outside the reached set, take a path
from the starting vertex to an unreached one and use its first edge
leaving the reached set. Adjoin that edge and its new endpoint.
Each step adds one new vertex attached by one edge, so it creates
no cycle. After one fewer additions than the number of vertices,
the selected edges span a connected acyclic graph, which is a
spanning tree. Conversely, a spanning tree supplies paths between
all vertices. A singleton uses the empty edge set. This proves the
equivalence with the source definition.

If \(C_\alpha\cap H\) has at least two elements, every pair of
them has the bidirected edge supplied by this same latent input.
They therefore lie in one district. A one-element intersection lies
in its unique district as well. Every input meeting \(H\) thus
belongs to exactly one displayed group \(\Lambda_D\); an input
meeting none belongs to none. The groups are disjoint and exhaustive
over the inputs that matter for the kernel.
\end{proof}
\end{annexitem}

\begin{annexitem}{annex:p16.latent-ancestral-reduction}{thm:p16-latent-ancestral-reduction}{Reverse-order summation at a fixed latent tuple}
\begin{proof}
First, closure under observed parents within \(H\) is equivalent
to closure under all directed ancestors there: repeated parent steps
give every directed ancestor, and a parent is itself an ancestor.
Use a topological order of \(G[H]\), and sum the indicator-product
formula for \(Q[H]\) over \(H\setminus A\) in reverse order,
holding the latent tuple \(u\) fixed during these summations.

When a coordinate \(j\in H\setminus A\) is reached, every
child still in \(H\setminus A\) has already been eliminated.
There is no child of \(j\) in \(A\), because that would make
\(j\) a parent in \(H\) of an \(A\) vertex. Consequently
the only remaining indicator using \(v_j\) is the one for its
own equation. Exactly one value of \(v_j\) satisfies that total
deterministic equation for the fixed parent inputs and \(u\).
Summing its indicator over \(S_j\) gives one and deletes it.
Induction along the reverse order leaves exactly the product of
equation indicators indexed by \(A\).

All sums are finite, so their order can be exchanged with the latent
sum. The resulting expression is the proved formula for \(Q[A]\).
No coordinate in \(H\setminus A\) is a parent of an \(A\)
vertex, so those extra outside assignments are unused in its
expression. This proves both the equality and its stated parameter
independence, including empty \(A\) or empty \(H\).
\end{proof}
\end{annexitem}

\begin{annexitem}{annex:p16.latent-target-reduction}{thm:p16-latent-target-reduction}{Target reduction before and after fixed actions}
\begin{proof}
With \(H=N\setminus X\), the law of the nonintervened observed
coordinates after \(X=x\) is \(Q[H]\). Thus
\[
 P_x(y)=\sum_{v_{H\setminus Y}}Q[H](v).
\]
The set \(A=\operatorname{an}_{G[H]}(Y)\) is ancestral in
\(G[H]\): appending a parent's edge to a path ending in \(Y\)
again gives such a path. Apply
\cref{thm:p16-latent-ancestral-reduction} to sum first over
\(H\setminus A\). This gives the displayed marginal of \(Q[A]\)
and proves independence of the arbitrary additional fixed values.
By the definition of \(Q[A]\), that same expression is the
response law under the joint intervention on \(X\) and
\(H\setminus A\), establishing the additional-action assertion.

For the reduction before intervention, \(A_0=\operatorname{an}_G(Y)\)
has no observed parent outside itself. The ancestral lemma with
\(H=N\) gives \(Q[A_0](v)=p_{A_0}(v_{A_0})\). Its retained equations
use only observed coordinates within \(A_0\) and their original
latent inputs. Evaluate them in a topological order, both in the
original model with intervention \(X=x\) and in this restricted
model with intervention \(X\cap A_0=x_{X\cap A_0}\). Their
retained and replaced equations agree, and so do their latent tuple
values. Induction gives equality of their solutions on \(A_0\),
including \(Y\), at every latent tuple. Their common latent law
then gives equality of the target distributions. No conditioning
event is introduced in either reduction.
\end{proof}
\end{annexitem}

\begin{annexitem}{annex:p16.latent-district-decomposition}{thm:p16-latent-district-decomposition}{Factoring the latent sum by its actual independent groups}
\begin{proof}
For fixed full \(v\), abbreviate the retained equation indicator by
\[
 \delta_i(u;v)=
 \mathbf1_{\{v_i=f_i(v_{\operatorname{pa}(i)},u_{\Lambda_i})\}}.
\]
Use the disjoint latent groups \(\Lambda_D\) from the partition
lemma. An indicator with \(i\in D\) uses only latent inputs in
\(\Lambda_D\): any input used there has child \(i\), and all
its other children within \(H\) belong to the same district.
An input with no child in \(H\) occurs in no indicator and its
probability mass sums to one. Finite distributivity of the remaining
independent latent sums therefore gives
\[
 Q[H](v)=\prod_{D\in\mathcal D_H}
 \left[
 \sum_{u_{\Lambda_D}}
      \prod_{\alpha\in\Lambda_D}r_\alpha(u_\alpha)
      \prod_{i\in D}\delta_i(u;v)
 \right].
\]
The latent inputs used by the equations in \(D\) are exactly
\(\Lambda_D\). Indeed any input with a child in \(D\) meets
\(H\), and its whole intersection with \(H\) lies in \(D\).
All other inputs sum to one in the full formula for \(Q[D]\).
Hence the bracketed expression is exactly \(Q[D](v)\), proving
the identity. All terms are nonnegative finite sums and products;
the proof uses no cancellation or division and includes zero values.

The normalization of each \(Q[D]\) is part of its kernel semantics,
with all other observed coordinates treated as fixed arguments.
In particular, an observed parent in another district is retained
as such an argument. This explains why the product does not assert
independence of the observed district tuples.
\end{proof}
\end{annexitem}

\begin{annexitem}{annex:p16.latent-prefix-chain}{thm:p16-latent-prefix-chain}{Prefix summation and positive conditional probabilities}
\begin{proof}
An observed parent within \(H\) of a vertex in \(H_j\) precedes
that vertex in the chosen order, and hence also belongs to \(H_j\).
Thus \(H_j\) is ancestral within \(H\), and the ancestral
reduction theorem gives \(\rho_j=Q[H_j]\). The marginal-sum
definition shows that \(\rho_j\) does not depend on the later
coordinates of \(H\).

For a fixed earlier prefix and outside parameters, summing \(\rho_j\)
over the next value gives \(\rho_{j-1}\). If the latter is
positive, division proves that \(\theta_j\) is a normalized
conditional mass. For any fixed full assignment, nonnegativity of
the summands gives
\(\rho_0\geq\rho_1\geq\cdots\geq\rho_k\geq0\), because
each prefix mass contains the mass at its specified next value.
Consequently \(\rho_{k-1}>0\) makes every denominator in the
product positive. Cancelling them in the finite product yields
\(\rho_k/\rho_0=Q[H]\). A zero \(\rho_k\) merely makes the
last factor zero; it is never a denominator in this calculation.
\end{proof}
\end{annexitem}

\begin{annexitem}{annex:p16.latent-district-recovery}{thm:p16-latent-district-recovery}{Grouping the conditional factors by district}
\begin{proof}
Fix the outside parameters and the full assignment \(v\). Keep
the district partition of the full current set \(H\), even if
a prefix splits one of those districts. For \(D\in\mathcal D_H\),
define
\[
 g_{D,j}(v)=\sum_{u_{\Lambda_D}}
       \prod_{\alpha\in\Lambda_D}r_\alpha(u_\alpha)
       \prod_{i\in D\cap H_j}\delta_i(u;v),
 \qquad 0\leq j\leq k,
\]
using the same equation indicators as in the district product proof.
The root groups are the groups for \(H\); roots not yet used
by a prefix have their normalized masses summed out. Applying finite
distributivity to the formula for \(Q[H_j]\) gives
\[
 \rho_j(v)=\prod_{D\in\mathcal D_H}g_{D,j}(v).
\]
Also \(g_{D,0}=1\), \(g_{D,k}=Q[D](v)\), and
\(g_{D,j}=g_{D,j-1}\) whenever \(h_j\notin D\), because
the index set of its equation indicators has then not changed.

If \(\rho_{j-1}>0\), each nonnegative factor \(g_{D,j-1}\)
is positive. Only the factor for the district \(D_j\) containing
\(h_j\) changes at this step. Cancellation of all other positive
factors therefore gives
\[
 \frac{\rho_j}{\rho_{j-1}}
      =\frac{g_{D_j,j}}{g_{D_j,j-1}}.
\]
Fix a district \(D\) and let \(m=m(D)\). Its assumed
\(\rho_{m-1}>0\), together with monotonicity of prefix masses,
makes every ratio at an index \(j\leq m\) well defined.
Multiply the displayed identity only over those indices belonging
to \(D\). The intervening unchanged values of \(g_{D,j}\)
make the product telescope from \(g_{D,0}=1\) to \(g_{D,m}\).
There are no later indices of \(D\), so
\(g_{D,m}=g_{D,k}=Q[D](v)\). Its final value may be zero;
all denominators used in telescoping are still positive. This proves
the stated identity on its exact sufficient domain.

If \(Q[H]\) is positive at every assignment for the fixed
outside parameters, every prefix contains a positive extension,
so the domain condition holds for each district and assignment.
Thus an available full kernel determines all its district kernels
on that parameter slice. For \(H=N\), \(Q[N]=p\) and its
prefix ratios are exactly the ordinary observed topological
conditionals. If two models in the specified class share this
strictly positive observed law, those ratios and all recovered
district factors agree. This is the asserted identification claim.
It does not assign values to a conditional ratio outside its support.
\end{proof}
\end{annexitem}

\begin{annexitem}{annex:p16.latent-foundation-calculation}{thm:p16-latent-foundation-calculation}{Every kernel, prefix and ancestral sum in the example}
\begin{proof}
The three root masses form the product
\(\frac12\alpha(e_A)\beta(e_B)\), summing to one. The observed
directed order is \(A,B,C\). The shared input \(U\) creates
the bidirected edge \(A\leftrightarrow C\), while the private
inputs create no bidirected edges. Thus the districts are exactly
\(\{A,C\}\) and \(\{B\}\), with latent groups
\(\{U,E_A\}\) and \(\{E_B\}\), respectively.

Fixing the other observed coordinates leaves the equation at \(B\)
as \(b=a\oplus E_B\). Its unique compatible disturbance is
\(E_B=b\oplus a\), proving \(Q[\{B\}]=\beta(b\oplus a)\).
Fixing \(B=b\) instead retains the equations for \(A,C\).
Summing out the private input for \(A\) gives
\[
 Q[\{A,C\}](a,b,c)=
 \tfrac12\sum_{u=0}^1\alpha(a\oplus u)
             \mathbf1_{\{c=\mathbf1_{\{b=1\text{ or }u=1\}}\}}.
\]
When \(b=0\), the response condition is \(u=c\), giving
\(\frac12\alpha(a\oplus c)\). When \(b=1\), it is
\(c=1\); the sum over \(u\) of \(\alpha(a\oplus u)\)
equals one, giving \(\frac12\mathbf1_{\{c=1\}}\).
These two kernels each sum to one over their own output coordinates
for any fixed outside arguments.

Multiplying by \(\beta(b\oplus a)\), as justified by the
district decomposition, yields for \((a,b)=(0,0)\) the two
masses \((3/4)(2/5)=12/40\) and \((3/4)(1/10)=3/40\).
For \((0,1)\), they are zero and \((1/4)(1/2)=5/40\).
For \((1,0)\), they are \((1/4)(1/10)=1/40\) and
\((1/4)(2/5)=4/40\). For \((1,1)\), they are zero and
\((3/4)(1/2)=15/40\). This proves all eight table entries,
which sum to one.

Summing the table gives \(p(A=a)=1/2\) and
\(p(A=a,B=b)=\frac12\beta(b\oplus a)>0\). Consequently
\(\theta_1=1/2\), \(\theta_2=\beta(b\oplus a)\), and
\(\theta_3=p(a,b,c)/[\frac12\beta(b\oplus a)]\).
The district with positions one and three is recovered as
\(\theta_1\theta_3=Q[\{A,C\}]\), and the other as
\(\theta_2=Q[\{B\}]\). When \(b=1,c=0\), only the final
numerator is zero; every required denominator remains positive.
Thus full positive joint support is not necessary in this example.

The induced graph on \(\{A,C\}\) has no directed edge, so
both singleton subsets are ancestral there. Summing the displayed
kernel over \(a\) gives \(Q[\{C\}](c;b)=1/2\) for either
\(c\) when \(b=0\), and gives point mass at \(c=1\) when
\(b=1\). Summing it over \(c\) gives \(Q[\{A\}]=1/2\)
for either \(a\), independent of the other fixed values.
In contrast, in the full observed graph the ancestral prefix
\(\{A,B\}\) has kernel
\(Q[\{A,B\}]=\frac12\beta(b\oplus a)\). Its induced graph
has two singleton districts; the shared input \(U\) now has
only its child \(A\) in the retained observed set. This checks
the restriction of latent child sets as well as the ancestral sums.
\end{proof}
\end{annexitem}

\begin{annexitem}{annex:p16.latent-foundation-limits}{thm:p16-latent-foundation-limits}{Three explicit limits of the graphical and ratio summaries}
\begin{proof}
For the first pair of models, take binary observed value sets and
no nonconstant latent input. Both evaluate observationally to
\((A,B)=(0,0)\). Their total equations at \(B\) are respectively
the constant zero and the parent input \(A\), on the same allowed
graph \(A\to B\). After fixing \(A=1\), their response-one
kernels are zero and one. In topological order \(A,B\), the
district \(\{B\}\) ends at position two and its required
denominator is \(\rho_1(A=1)=p(A=1)=0\). Both the numerator
and denominator of the proposed observed conditional ratio are zero.
The structural kernels are still defined, but this observed law
does not determine the unobserved-input response value.

For the second pair, the four input tuples \((u,t)\) have mass
\(1/4\), proving independence of the two fair latent roots.
The declared observed graph has \(A\to Y\) and
\(A\leftrightarrow Y\). In the first model, observationally
\(A=U\) makes \(U\oplus A=0\), so \(Y=T\). The second
model has \(Y=T\) by its equation. Thus both full observed
tables have mass \(1/4\) at every \((a,y)\), and are strictly
positive. Under \(A=a\), the first model has response zero
exactly when \(T=0\) and \(U=a\). That event has probability
\(1/4\), so its response-one kernel is \(3/4\). The second
model retains response-one probability \(1/2\).

The singleton \(\{Y\}\) is connected by its length-zero path,
but the full graph has the larger district \(\{A,Y\}\).
Recovery of the latter gives the common observed joint law; it
does not identify the former. This is a positive-law example of
the source's bow-arc obstruction
\cite[Theorem~2, p.~1221]{ShpitserPearl2006JointIdentification},
with its extra random response branch written explicitly. The
constant/ignored arguments in the second model retain their allowed
input declarations, as permitted by the model class. No faithfulness
or minimal dependence graph is assumed.

For the last comparison, the single fair input gives \((A,B,C)\)
equal to \((0,0,0)\) or \((1,1,1)\) with mass \(1/2\)
each. Its child set contains all three vertices and therefore
creates all three bidirected edges. For the pair-input model let
the independent fair bits be \(R_{AB},R_{AC},R_{BC}\) and put
\[
 A=R_{AB}\oplus R_{AC},\qquad
 B=R_{AB}\oplus R_{BC},\qquad
 C=R_{AC}\oplus R_{BC}.
\]
Each root is incident to exactly its named pair, producing the
same declared mixed graph. The output parity is zero, since each
input occurs twice in \(A\oplus B\oplus C\). For any prescribed
even-parity output \((a,b,c)\), choose \(R_{AB}=t\) freely;
then \(R_{AC}=t\oplus a\) and \(R_{BC}=t\oplus b\)
are forced, and their last output equals \(a\oplus b=c\).
There are exactly two compatible inputs among the eight equally
likely tuples. Each of \(000,011,101,110\) therefore has mass
\(1/4\), and all odd-parity outputs have mass zero. These
probabilities differ from the common-input construction. This
calculation proves the stated failure of that replacement, without
a claim about every alternative pairwise construction.
\end{proof}
\end{annexitem}
